\section*{Chapter \ref{ch:s1:positioning-and-sentiment} --- Positioning and Sentiment}

\begin{solution}{exo:s1:positioning-and-sentiment:1}
From Friday's close (after the 3:30 pm release), or Monday's open to be safe; later in holiday weeks. Using it on Tuesday gives the backtest three days of
information nobody had, and credits it with the signal's freshest, most valuable part.
\end{solution}

\begin{solution}{exo:s1:positioning-and-sentiment:2}
$\phi = 2^{-1/20}$: $\phi^3 = 0.901$, $\phi^{10} = 0.707$, $\phi^{21} = 0.483$. The measured ICs fall a little faster (to 86\%, 66\% and 38\% of the
immediate one) because a report is used for five days after it becomes known.
\end{solution}

\begin{solution}{exo:s1:positioning-and-sentiment:3}
Every futures contract has a long and a short. The report's categories cover almost all open interest, and in corn the two largest net positions are
hedgers' (producers and merchants, net short) and speculators' (managed money): when one adds, the other takes the other side.
\end{solution}

\begin{solution}{exo:s1:positioning-and-sentiment:4}
About $1.96/\sqrt{113} = 0.184$. Corn's 0.098 is about half of that ($t = 1.04$).
\end{solution}

\begin{solution}{exo:s1:positioning-and-sentiment:5}
Part of the speculators' position is trend-chasing, and \texttt{firm.synthfut}'s planted trends have a half-life of a year: that part of the signal
predicts returns for longer than the \omterm{def:s1:positioning-and-sentiment:hedging}{hedging pressure}, whose half-life is 20 days.
\end{solution}

\begin{solution}{exo:s1:positioning-and-sentiment:6}
In the model speculators are paid for absorbing hedging, so large speculative longs come with high expected returns: fading them loses. A contrarian rule
needs speculators to be the ones paying: positions that push prices away from value and later unwind, a crowding mechanism the model does not have.
\end{solution}

\begin{solution}{exo:s1:positioning-and-sentiment:7}
The IC falls from 0.032 known at once to 0.019 known on Friday: the lag costs 42\% of the signal, against 14\% with a half-life of twenty days. The
shorter the pressure lasts, the more a three-day delay costs.
\end{solution}

\begin{solution}{exo:s1:positioning-and-sentiment:8}
Survey values dated by their reference period, not by their publication, give the backtest information before it was public; indices built with
principal components over the whole sample use later data to weight earlier values; and the proxies were chosen after the fact. Use publication dates
and an index rebuilt with data available at each date.
\end{solution}

\begin{solution}{pb:s1:positioning-and-sentiment:1}
\begin{enumerate}
\item Generally Friday at 3:30 pm Eastern Time, for the preceding Tuesday.
\item Hedgers' net position as a share of open interest; a signal from reported positions as known at the time.
\item Returns in currency and agricultural futures vary with hedgers' net holdings; own and within-group \omterm{def:s1:positioning-and-sentiment:hedging}{hedging pressure} affect returns in 20 markets.
\item As mirror images: a correlation of $-0.97$.
\item Managed money's 52-week z-score at each month-end against the maize price's change over the month after next, so that the signal's month does not
overlap the return.
\item 0.098, $t = 1.04$.
\item $+4.6\%$ over three months after record longs, $+0.2\%$ after record shorts, $-0.9\%$ otherwise.
\item Nothing either way.
\item An AR(1) with a half-life of 20 days earning 0.3 volatilities a year per standard deviation; speculators hold it plus trend-chasing and noise.
\item 0.044, 0.038, 0.029 and 0.017 at lags of 0, 3, 10 and 21 days.
\item Three days keep $\phi^3 = 0.90$ of the pressure; a report used for five days loses a little more: 14\%.
\item Speculators are paid in the model, so their extremes come with high expected returns.
\item \textbf{Named result.} The hedgers' \omterm{def:s1:positioning-and-sentiment:hedging}{positioning signal} has a rank IC of 0.044 with the next five days' returns when known at once and 0.038 when
published three days later (14\% less), falling to 0.017 at 21 days; the speculators' signal, 0.043 and 0.039.
\item A summary of investors' optimism from surveys or market proxies; low sentiment is followed by relatively high returns on hard-to-value stocks.
\item As a slow conditioning variable, saying which stocks are most mispriced, not as a weekly trading signal.
\item The positions of those paying for convexity, daily and by strike.
\item Reports only from their publication date, holiday shifts included, and costs.
\item The extreme-positioning contrarian.
\item Both pay whoever carries risk that others want to lay off: carry through the curve, \omterm{def:s1:positioning-and-sentiment:hedging}{hedging pressure} through positions.
\item Hedgers who pay to lay off price risk.
\end{enumerate}
\end{solution}

\begin{solution}{iq:s1:positioning-and-sentiment:1}
Open interest and the long and short positions of reportable trader categories (producers and merchants, swap dealers, managed money, others), weekly
for Tuesday; it cannot tell you positions within the week, positions in other markets or over the counter, or why a category trades.
\end{solution}

\begin{solution}{iq:s1:positioning-and-sentiment:2}
Hedgers pay to transfer price risk: producers selling futures below the expected spot price leave a premium for the speculators who buy them.
\end{solution}

\begin{solution}{iq:s1:positioning-and-sentiment:3}
A record long alone is not a signal: check whether it is absorbing hedging (producers' positions, curve shape, inventories) or crowding on a trend;
check the evidence for that market; size any fade small, because extremes can persist and in some data are followed by further gains.
\end{solution}

\begin{solution}{iq:s1:positioning-and-sentiment:4}
Each report with its as-of date (Tuesday) and its publication timestamp, from the CFTC's release calendar; revisions kept as new versions; queries
always by publication time.
\end{solution}

\begin{solution}{iq:s1:positioning-and-sentiment:5}
Choose proxies with a reason (surveys, issuance, flows, new-issue returns), orthogonalise them to macroeconomic conditions, combine them with weights
estimated only on past data, and test the index on the cross-section of hard-to-value stocks out of sample.
\end{solution}

\begin{solution}{iq:s1:positioning-and-sentiment:6}
With $r_{t+1} = a s_t + e_{t+1}$, $r_{t+k} = a s_{t+k-1} + e_{t+k}$ and $s_{t+k-1} = \phi^{k-1}s_t + \text{independent terms}$, so the correlation of $s_t$
with $r_{t+k}$ is $\rho\phi^{k-1}$. A report used on days $t + \ell, \dots, t + \ell + 4$ has average correlation $\rho\phi^{\ell}(1 - \phi^5)/(5(1 - \phi))$.
\end{solution}
